% E01 report. Built by `make e01`; numbers come from build/E01/results.tex.
\documentclass[zihao=-4, linespread=1.5, fontset=none]{ctexart}
\usepackage{labreport}

\graphicspath{{../../build/E01/}}
\input{../../build/E01/results.tex}

\newcommand{\src}{../../firmware/audio_loopback/dsp}

\title{实验 E01\quad 数字滤波器设计与实现}
\author{}
\date{}

\begin{document}
\maketitle

\section{实验内容}

在 ADSP-SC589 EZ-Board 上实现两个低通滤波器，采样率 $f_s=\SampleRate$~Hz，截止频率 $f_c=\CutoffHz$~Hz：
\begin{itemize}
  \item FIR 滤波器，\FirOrder{} 阶（\FirTaps{} 个抽头）；
  \item IIR 滤波器，\IirOrder{} 阶，由 \IirSections{} 个二阶节级联。
\end{itemize}
通过分析滤波器的输入与输出，比较二者的单位冲激响应、幅频特性、相频特性、群时延与计算量。

\section{系统框图}

板上程序把 ADC 采到的音频按块送入处理函数，再送给 DAC 播放（图~\ref{fig:system}）。
ADAU1979 输出 24 位补码样本，经 SPORT 与 DMA 以每块 \BlockFrames{} 帧的乒乓缓冲送进 SHARC 核。
\texttt{AudioProcess} 把每个声道解交织并转为 $[-1,1)$ 的浮点数，偶数声道走 FIR、奇数声道走 IIR，
再饱和、舍入回 24 位交给 ADAU1962。同一路输入接到一对相邻声道上，就能在两个输出上同时听到、测到两种滤波器。

\begin{figure}[H]
\centering
\begin{tikzpicture}[node distance=6mm and 7mm, font=\small]
  \node[block] (adc) {ADC\\ADAU1979};
  \node[block, right=of adc] (rx) {SPORT 接收\\DMA 乒乓};
  \node[block, right=of rx] (cvt) {解交织\\24 位 $\to$ float};
  \node[block, right=12mm of cvt, yshift=8mm] (fir) {FIR\\\FirTaps{} 抽头};
  \node[block, right=12mm of cvt, yshift=-8mm] (iir) {IIR\\\IirSections{} 个二阶节};
  \node[block, right=12mm of fir, yshift=-8mm] (back) {float $\to$ 24 位\\交织};
  \node[block, below=10mm of back] (tx) {SPORT 发送\\DMA 乒乓};
  \node[block, left=of tx] (dac) {DAC\\ADAU1962};
  \draw[sig] (adc) -- (rx);
  \draw[sig] (rx) -- (cvt);
  \draw[sig] (cvt.east) -- ++(4mm,0) |- node[pos=0.75, above] {\scriptsize 偶声道} (fir.west);
  \draw[sig] (cvt.east) -- ++(4mm,0) |- node[pos=0.75, below] {\scriptsize 奇声道} (iir.west);
  \draw[sig] (fir.east) -| ($(back.west)+(-4mm,0)$) -- (back.west);
  \draw[sig] (iir.east) -| ($(back.west)+(-4mm,0)$) -- (back.west);
  \draw[sig] (back) -- (tx);
  \draw[sig] (tx) -- (dac);
  \node[draw, dashed, rounded corners, inner sep=3mm, fit=(cvt) (fir) (iir) (back),
        label={[font=\small]above:\texttt{AudioProcess}（SHARC 核 1）}] {};
\end{tikzpicture}
\caption{板上音频通路}
\label{fig:system}
\end{figure}

\section{滤波器设计}

两组系数都由 \texttt{coursework/E01-fir-iir/design.py} 用 SciPy 计算，并写成 C 数组
\texttt{lowpass\_coeffs.c}，固件直接编译进去。

\subsection{FIR：窗函数法}

理想低通的冲激响应是 sinc 函数，截取以 $(N-1)/2$ 为中心的 $N=\FirTaps$ 点，再乘 \FirWindow{} 窗 $w[n]$：
\begin{equation}
  h[n] = w[n]\,\frac{2f_c}{f_s}\,\operatorname{sinc}\!\left(\frac{2f_c}{f_s}\Bigl(n-\frac{N-1}{2}\Bigr)\right),
  \qquad n=0,\dots,N-1,
\end{equation}
最后归一化使 $\sum h[n]=1$（直流增益为 1）。$f_c$ 是理想响应的半幅点（$-6$~dB），
这也是 MATLAB \texttt{fir1} 的定义；\FirTaps{} 抽头的实际滤波器在 $f_c$ 处为 \FirAtCutoff~dB。
$h[n]$ 关于中点偶对称，因此相位严格线性。

FIR 用直接型（横向）结构实现（图~\ref{fig:fir}），每个输出样本做 \FirMuls{} 次乘法、\FirAdds{} 次加法，
延时线保存 \FirStates{} 个样本。

\begin{figure}[H]
\centering
\begin{tikzpicture}[node distance=5mm and 9mm, font=\small]
  \node (x) {$x[n]$};
  \node[dot, right=6mm of x] (d0) {};
  \node[block, right=of d0, minimum width=2.6em] (z1) {$z^{-1}$};
  \node[dot, right=of z1] (d1) {};
  \node[block, right=of d1, minimum width=2.6em] (z2) {$z^{-1}$};
  \node[dot, right=of z2] (d2) {};
  \node[right=6mm of d2] (dots) {$\cdots$};
  \node[block, right=6mm of dots, minimum width=2.6em] (zn) {$z^{-1}$};
  \node[dot, right=of zn] (dn) {};
  \foreach \d/\k in {d0/0, d1/1, d2/2, dn/N-1} {
    \node[draw, isosceles triangle, shape border rotate=-90, inner sep=0pt,
          minimum height=6mm, minimum width=8mm, below=6mm of \d, font=\tiny] (g\k) {$h_{\k}$};
  }
  \node[sum, below=6mm of g0] (s0) {$+$};
  \node[sum, below=6mm of g1] (s1) {$+$};
  \node[sum, below=6mm of g2] (s2) {$+$};
  \node[sum, below=6mm of gN-1] (sn) {$+$};
  \node[right=8mm of sn] (y) {$y[n]$};
  \draw (x) -- (d0);
  \draw[sig] (d0) -- (z1); \draw (z1) -- (d1);
  \draw[sig] (d1) -- (z2); \draw (z2) -- (d2);
  \draw (d2) -- (dots); \draw[sig] (dots) -- (zn); \draw (zn) -- (dn);
  \draw[sig] (d0) -- (g0); \draw[sig] (g0) -- (s0);
  \draw[sig] (d1) -- (g1); \draw[sig] (g1) -- (s1);
  \draw[sig] (d2) -- (g2); \draw[sig] (g2) -- (s2);
  \draw[sig] (dn) -- (gN-1); \draw[sig] (gN-1) -- (sn);
  \draw[sig] (s0) -- (s1);
  \draw[sig] (s1) -- (s2);
  \draw (s2) -- node[fill=white, inner sep=1pt] {$\cdots$} (sn.west -| dots);
  \draw[sig] (sn.west -| dots) -- (sn);
  \draw[sig] (sn) -- (y);
\end{tikzpicture}
\caption{FIR 直接型结构，$N=\FirTaps$}
\label{fig:fir}
\end{figure}

\subsection{IIR：Butterworth 二阶节级联}

\IirOrder{} 阶 Butterworth 低通由双线性变换得到，$f_c$ 处为 $-3$~dB。直接型高阶 IIR 对系数量化很敏感，
因此按极点对拆成 \IirSections{} 个二阶节级联：
\begin{equation}
  H(z) = \prod_{k=1}^{\IirSections} \frac{b_{0k} + b_{1k}z^{-1} + b_{2k}z^{-2}}{1 + a_{1k}z^{-1} + a_{2k}z^{-2}} .
\end{equation}
SciPy 把总增益（约 $5\times10^{-8}$）全放在第一节；设计脚本把每一节各自缩放到直流增益为 1，
避免中间信号贴近浮点下限后再被后级放大。

每个二阶节用转置直接 II 型实现（图~\ref{fig:biquad}），只需 2 个状态变量：
\begin{equation}
  y = b_0x + s_1,\qquad s_1 \leftarrow b_1x - a_1y + s_2,\qquad s_2 \leftarrow b_2x - a_2y .
\end{equation}
整个级联每个输出样本做 \IirMuls{} 次乘法、\IirAdds{} 次加法，保存 \IirCoeffs{} 个系数和 \IirStates{} 个状态。

\begin{figure}[H]
\centering
\begin{tikzpicture}[font=\small, node distance=8mm and 12mm]
  \node (x) {$x$};
  \node[dot, right=8mm of x] (xd) {};
  \node[sum, right=26mm of xd] (s0) {$+$};
  \node[dot, right=10mm of s0] (yd) {};
  \node[right=8mm of yd] (y) {$y$};
  \node[block, below=7mm of s0, minimum width=2.6em] (z1) {$z^{-1}$};
  \node[sum, below=7mm of z1] (s1) {$+$};
  \node[block, below=7mm of s1, minimum width=2.6em] (z2) {$z^{-1}$};
  \node[sum, below=7mm of z2] (s2) {$+$};
  \node[left=13mm of s0, font=\scriptsize] (b0) {$b_0$};
  \node[left=13mm of s1, font=\scriptsize] (b1) {$b_1$};
  \node[left=13mm of s2, font=\scriptsize] (b2) {$b_2$};
  \node[right=13mm of s1, font=\scriptsize] (a1) {$-a_1$};
  \node[right=13mm of s2, font=\scriptsize] (a2) {$-a_2$};
  \draw (x) -- (xd);
  \draw[sig] (xd) -- (b0) ; \draw[sig] (b0) -- (s0);
  \draw (xd) |- (b1); \draw[sig] (b1) -- (s1);
  \draw (xd) |- (b2); \draw[sig] (b2) -- (s2);
  \draw (s0) -- (yd); \draw[sig] (yd) -- (y);
  \draw (yd) |- (a1); \draw[sig] (a1) -- (s1);
  \draw (yd) |- (a2); \draw[sig] (a2) -- (s2);
  \draw[sig] (s2) -- node[right, font=\scriptsize] {} (z2);
  \draw[sig] (z2) -- (s1);
  \draw[sig] (s1) -- (z1);
  \draw[sig] (z1) -- (s0);
  \node[right=1mm of z1, font=\scriptsize] {$s_1$};
  \node[right=1mm of z2, font=\scriptsize] {$s_2$};
\end{tikzpicture}
\caption{二阶节的转置直接 II 型结构}
\label{fig:biquad}
\end{figure}

\section{关键代码}

FIR 的延时线是长度 $N$ 的环形缓冲。为了不在乘累加里做取模，卷积按绕回点拆成两段循环：

\lstinputlisting[language=C]{\src/fir.c}

IIR 逐节计算，前一节的输出就是后一节的输入：

\lstinputlisting[language=C]{\src/biquad.c}

系数设计：

\lstinputlisting[language=Python, rangeprefix=, rangesuffix=, includerangemarker=true,
  linerange={def\ fir(fs):-return\ sos}]{design.py}

\section{结果与分析}

板上程序启动时向每个滤波器输入单位冲激 $\delta[n]$，把输出 $y[n]$ 逐点打印到串口
（FIR \FirTaps{} 点之后再多打同样多个零，IIR 打 \IirReportLen{} 点），并用周期计数器测量处理一块
\BlockFrames{} 帧所用的周期数。幅频、相频与群时延由设计系数计算；板上冲激响应与设计的冲激响应逐点吻合
（最大误差 FIR \FirIrErr、IIR \IirIrErr，即单精度浮点的舍入量级），所以板上滤波器的频域特性就是图中的设计曲线。

\subsection{单位冲激响应}

\begin{figure}[H]
\centering
\includegraphics[width=\linewidth]{impulse.pdf}
\caption{单位冲激响应。竖线为设计值，叉为板上 SHARC 的输出；IIR 只画前 \IirShown{} 点}
\label{fig:impulse}
\end{figure}

FIR 的冲激响应就是 \FirTaps{} 个系数本身，长度有限，第 \FirTaps{} 点之后严格为零，关于第 \FirGroupDelay{} 点对称。
IIR 的冲激响应无限长，从零缓慢升起，峰值只有 \IirIrPeak{}，之后以衰减振荡的形式拖尾：
极点靠近单位圆，能量释放得慢。FIR 的对称性对应线性相位，IIR 的不对称、长拖尾对应下面看到的非线性相位。

\subsection{幅频特性}

\begin{figure}[H]
\centering
\includegraphics[width=\linewidth]{magnitude.pdf}
\includegraphics[width=\linewidth]{magnitude_pass.pdf}
\caption{幅频特性：上为 $0$--$f_s/2$，下为通带与过渡带}
\label{fig:magnitude}
\end{figure}

\begin{table}[H]
\centering
\caption{各频率点的增益（dB）}
\label{tab:gain}
\begin{tabular}{ccc}
\toprule
频率 / kHz & FIR & IIR \\
\midrule
\GainRows
\bottomrule
\end{tabular}
\end{table}

\begin{itemize}
  \item \textbf{截止点。} IIR 在 \CutoffHz~Hz 处为 \IirAtCutoff~dB，符合 Butterworth 的 $-3$~dB 定义；
    FIR 在同一点为 \FirAtCutoff~dB，二者截止频率的定义不同。
  \item \textbf{通带。} Butterworth 通带最大平坦，直到接近 $f_c$ 都几乎没有衰减；
    FIR 的增益在通带内就开始缓慢下滑（表~\ref{tab:gain}）。
  \item \textbf{过渡带与阻带。} 同样以 $f_c$ 为界，\IirOrder{} 阶 IIR 的过渡带明显更陡（每倍频程约 48~dB）；
    \FirOrder{} 阶 FIR 的过渡带宽，阻带衰减受 \FirWindow{} 窗旁瓣的限制，起伏而不再随频率增加，
    而 IIR 的衰减随频率单调增加，到 $f_s/2$ 处趋于 $-\infty$（双线性变换把 $s=\infty$ 的零点映射到 $z=-1$）。
  \item 要让 FIR 达到与 IIR 相当的过渡带，阶数需要高得多，这是 IIR 在幅频选择性上的优势。
\end{itemize}

\subsection{相频特性与群时延}

\begin{figure}[H]
\centering
\includegraphics[width=\linewidth]{phase.pdf}
\includegraphics[width=\linewidth]{group_delay.pdf}
\caption{相频特性（展开后）与群时延，$0$--4~kHz}
\label{fig:phase}
\end{figure}

\begin{itemize}
  \item \textbf{FIR 线性相位。} 系数偶对称，$\varphi(\omega)=-\frac{N-1}{2}\omega$，相频曲线是直线，
    群时延恒为 $\frac{N-1}{2}=\FirGroupDelay$ 个样本（\FirGroupDelayMs~ms），所有频率分量延时相同，波形不失真。
  \item \textbf{IIR 非线性相位。} 相位在截止频率附近弯折，群时延从直流处的 \IirGroupDelayDc{} 个样本
    上升到 \IirGroupDelayPeakHz~Hz 处的峰值 \IirGroupDelayPeak{} 个样本（\IirGroupDelayPeakMs~ms）。
    通带内不同频率的延时不同，含有接近截止频率成分的信号（如方波、瞬态）会出现相位失真和振铃。
  \item 就延时本身而言，本设计的 IIR 在通带内的群时延反而大于 FIR；IIR 的优势在计算量与选择性，不在延时。
\end{itemize}

\subsection{计算量}

\begin{table}[H]
\centering
\caption{计算量比较（每个输出样本）}
\label{tab:cost}
\begin{tabular}{lcc}
\toprule
 & FIR（\FirOrder{} 阶） & IIR（\IirOrder{} 阶） \\
\midrule
乘法次数 & \FirMuls & \IirMuls \\
加法次数 & \FirAdds & \IirAdds \\
系数个数 & \FirTaps & \IirCoeffs \\
状态（延时单元）个数 & \FirStates & \IirStates \\
实测周期数（\BlockFrames{} 帧一块） & \FirCycles & \IirCycles \\
实测周期 / 样本 & \FirPerSample & \IirPerSample \\
\bottomrule
\end{tabular}
\end{table}

周期数在 CCLK $=\CclkMHz$~MHz 的 SHARC 核上测得，程序为 Debug 配置（未开编译优化）。每个样本的周期预算为
$\CclkMHz\,\text{MHz}/f_s=\CyclesBudget$；处理 4 个声道（2 路 FIR、2 路 IIR）连同格式转换的整个
\texttt{AudioProcess} 每帧用 \AudioPerFrame{} 个周期，占预算的 \AudioLoad\%。

\begin{itemize}
  \item 按运算次数，FIR 的乘法是 IIR 的 \FirOverIirOps{} 倍；实测周期 FIR 是 IIR 的 \FirOverIir{} 倍。
    差距更大，是因为未优化的代码每做一次乘加都有一套循环开销，FIR 内层循环每个样本跑 \FirTaps{} 次，
    IIR 只跑 \IirSections{} 次。
  \item 要达到相近的幅频选择性，FIR 需要的阶数远高于 IIR，计算量和存储的差距还会进一步拉大。
  \item FIR 没有反馈，系数量化和舍入误差不会累积，总是稳定的；IIR 有反馈，高阶时必须拆成二阶节，
    极点靠近单位圆时对系数精度更敏感。
\end{itemize}

\section{结论}

\begin{table}[H]
\centering
\caption{两种滤波器的比较}
\begin{tabular}{lll}
\toprule
 & FIR（\FirOrder{} 阶，\FirWindow{} 窗） & IIR（\IirOrder{} 阶 Butterworth） \\
\midrule
冲激响应 & 有限长，偶对称 & 无限长，衰减振荡 \\
\CutoffHz~Hz 处增益 & \FirAtCutoff~dB & \IirAtCutoff~dB \\
过渡带 & 宽 & 窄 \\
相位 & 线性 & 非线性 \\
群时延 / 样本 & 恒为 \FirGroupDelay & \IirGroupDelayDc{}--\IirGroupDelayPeak \\
乘法 / 样本 & \FirMuls & \IirMuls \\
稳定性 & 恒稳定 & 取决于极点 \\
\bottomrule
\end{tabular}
\end{table}

需要波形不失真（线性相位）时选 FIR；更看重以较少运算换取陡峭的过渡带时选 IIR。

\end{document}
